\section{Introduction}
\label{sec:intro}

Estimates pervade modern programs. GPS readings, sensor streams, machine-learning
predictions, and survey data all arrive carrying error. Discrete types
(\texttt{float}, \texttt{int}, \texttt{bool}) encourage developers to treat
these estimates as facts, and Bornholt et al.~\cite{bornholt2014uncertain}
characterised three resulting bug classes: using estimates as facts ignores
random error, computation compounds that error, and Boolean questions on
probabilistic data create false positives and false negatives. Their
\texttt{Uncertain<T>} type addresses all three by lifting random variables into
the type system, propagating distributions through arithmetic, and discharging
conditionals via Wald's sequential probability ratio test~\cite{wald1945sprt}.

\texttt{Uncertain<T>} assumes the value exists. The only uncertainty it models
is in the measurement. Many real workloads violate this assumption. A patient
in a headache trial reports pain reduction on some days and not others;
when a report arrives, the reduction itself is noisy. A temperature sensor on
an embedded device occasionally drops a sample;
when a sample arrives, the reading carries Gaussian noise. A vehicle entering
a tunnel loses its GPS fix for fifteen seconds;
when fixes resume, each carries a horizontal-accuracy radius.
Three failure modes appear in practice when programmers handle such data with
existing tools.

\paragraph{Phantom presence.} The simplest workaround is to assume the value
always arrives and assign a wide distribution to the missing case. This
silently injects a fabricated number into downstream arithmetic. The result
looks uncertain, in the type-system sense, but the uncertainty has no
empirical grounding.

\paragraph{Hard-gated drop.} A more careful programmer wraps the type as
\texttt{Option<Uncertain<T>>} and calls \texttt{is\_some} before using the
value. This collapses the probabilistic presence into a deterministic Boolean
at the boundary, recreating the very class of false-positive and false-negative
bugs that motivated \texttt{Uncertain<T>}.

\paragraph{Imputation laundering.} Statistical practice replaces missing
values with means, zeros, or model-based imputations. The imputed quantity then
flows through computation indistinguishable from real data. Inside a
probabilistic type system the laundering is invisible: the type says
``uncertain value,'' the underlying distribution is whatever the imputation
chose, and no flag records that the number was synthesised rather than
measured.

\medskip

The statistical literature names this problem cleanly. Rubin~\cite{rubin1976missing}
distinguishes missing completely at random (MCAR), missing at random (MAR), and
missing not at random (MNAR), and a large body of work explores estimation
under each regime~\cite{littlerubin2019}. What is missing is a programming
abstraction that lets the developer state, at the level of a single value,
which assumption applies and propagates that assumption through subsequent
computation.

This paper proposes such an abstraction. \texttt{MaybeUncertain<T>} factors
the two sources of uncertainty into a pair of independent random variables:
$\mathit{is\_present}$, an \texttt{Uncertain<bool>} drawn from a Bernoulli
distribution that governs whether the value exists at all, and
$\mathit{value}$, an \texttt{Uncertain<T>} describing the value's distribution
conditional on presence. Arithmetic operators propagate both axes; presence
combines under AND. A statistically-gated operator,
\texttt{lift\_to\_uncertain}, returns a plain \texttt{Uncertain<T>} only when
an SPRT on the presence axis clears a confidence threshold; otherwise it
returns a \texttt{PresenceError}.

This is an extension, not a replacement. \texttt{MaybeUncertain<T>} reuses
Bornholt's Bayesian-network sampling semantics, the SPRT machinery at
conditionals, and the global sample cache. The new construction is one
additional Bernoulli node attached to each leaf, plus AND-propagation on
arithmetic, plus the gating operator. Per-sample overhead in our benchmarks
falls between $1.6\times$ and $2\times$; the SPRT gate costs a few hundred
nanoseconds when the evidence is decisive and a few microseconds when the true
presence probability is close to the threshold.

\paragraph{Contributions.} This paper makes four contributions.

\begin{enumerate}
  \item A taxonomy of three sparse-data uncertainty bugs (phantom presence,
        hard-gated drop, imputation laundering) that the original
        \texttt{Uncertain<T>} taxonomy does not catch.
  \item The \texttt{MaybeUncertain<T>} type, formalised as a product of two
        independent \texttt{Uncertain} random variables, with operator
        semantics that propagate presence under AND.
  \item The \texttt{lift\_to\_uncertain} gating operator, which uses SPRT on
        the presence axis to expose the inner value type with a stated
        confidence level.
  \item An open-source Rust implementation in the \texttt{deep\_causality\_uncertain}
        crate of the DeepCausality project~\cite{hansen_deepcausality}, with
        three worked examples (clinical trial, intermittent sensor, GPS dropouts)
        and a Criterion benchmark suite quantifying the overhead.
\end{enumerate}

\Cref{sec:background} surveys the prior work that
\texttt{MaybeUncertain<T>} extends and the adjacent communities whose
constructions partially overlap with ours. \Cref{sec:type}
defines the type, its constructors, and its algebra. \Cref{sec:semantics}
gives the sampling semantics, the Bayesian-network embedding, and the SPRT
hypothesis at the gate. \Cref{sec:implementation} describes the Rust API and
the reuse of existing infrastructure. \Cref{sec:evaluation} reports
benchmark results, and \Cref{sec:examples} presents the three worked examples.
\Cref{sec:discussion} addresses limitations honestly, and
\Cref{sec:conclusion} concludes.
